% ------------------------------------------------------------------------------------
%                       SODDISFACIBILITÀ BOOLEANA
% ------------------------------------------------------------------------------------

\chapter{Soddisfacibilità Booleana}

Abbiamo già incontrato brevemente i \textbf{problemi SAT} 
(\textit{Boolean Satisfiability}) come esempio di classe di problemi \textbf{facili da 
modellare} con il paradigma MIP, ma \textbf{difficili da risolvere} tramite rilassamenti 
continui.

In questo capitolo approfondiamo i \textbf{fondamenti teorici} della 
\textbf{soddisfacibilità booleana} e i principali \textbf{algoritmi risolutivi} 
dedicati a questa classe di problemi.

\section{Logica delle proposizioni}

La \textbf{logica delle proposizioni} è lo studio di \textbf{formule} la cui verità è 
determinata dal \textbf{valore di verità delle loro parti}. Tali formule sono ottenute 
combinando formule più piccole tramite \textbf{connettivi logici}, quali \textit{and} ($\land$), 
\textit{or} ($\lor$) e \textit{not} ($\neg$).

I componenti di base sono detti \textbf{proposizioni atomiche}. Una formula booleana 
ottenuta combinando le proposizioni atomiche $x_1, \ldots, x_n$ può pertanto essere 
vista come una \textbf{funzione booleana}:
\[
    f(x_1, \ldots, x_n) : \{T, F\}^n \to \{T, F\}
\]
Dove con $T$ e $F$ intendiamo i valori di verità \textit{True} e \textit{False} 
(che, come abbiamo già visto, possiamo mappare rispettivamente nei valori $1$ e $0$).

\begin{definition}[FORMULA LOGICA]
    Più formalmente, le \textbf{formule logiche} possono essere definite 
    \textbf{ricorsivamente} come:
    \begin{itemize}
        \item Una \textbf{formula \textit{vuota}} (per definizione \textit{falsa}).
        \item Una \textbf{proposizione atomica} $x_j$, ovvero una formula \textit{"di base"} 
            che può assumere valore vero o falso.
        \item Un'\textbf{espressione \textit{composta}} del tipo $\neg(F)$, $(F) \lor (G)$ e 
            $(F) \land (G)$, dove $F$ e $G$ sono a loro volta formule logiche.
    \end{itemize}
\end{definition}

Chiaramente, il valore di \textbf{verità di una formula composta dipende} dal valore di 
\textbf{verità delle formule costituenti} e dalle \textbf{tabelle di verità} associate 
ai \textbf{connettivi utilizzati}.

Oltre ai connettivi di base ($\land$, $\lor$, $\neg$), sono spesso usati i 
connettivi di \textbf{implicazione} ($\to$) e di \textbf{equivalenza} ($\equiv$), 
che risultano tuttavia \textbf{esprimibili tramite i connettivi base}:

\begin{remark}[CONNETTIVI DERIVATI]
    \[
        \begin{aligned}
            F \to G &\equiv \neg F \lor G \\
            (F \equiv G) &\equiv (F \to G) \land (G \to F)
        \end{aligned}
    \]
\end{remark}

\section{Il problema SAT}

Per poter trattare algoritmicamente le formule logiche, conviene restringersi a una 
\textbf{struttura sintattica standardizzata}. Introduciamo quindi i concetti di 
\textit{literal}, \textit{clausola} e \textit{forma congiuntiva normale}.

\begin{definition}[LITERAL E CLAUSOLA]
    Si definisce \textbf{literal} una \textbf{proposizione atomica} $x_j$ oppure la 
    \textbf{sua negazione} $\neg x_j$.

    Si definisce poi \textbf{clausola} (\textit{clause}) una \textbf{disgiunzione} ($\lor$) 
    di \textit{literals}.
\end{definition}

\begin{definition}[FORMA CONGIUNTIVA NORMALE (CNF)]
    Una formula logica si dice in \textbf{forma congiuntiva normale} 
    (\textit{Conjunctive Normal Form}, \textbf{CNF}) se è espressa come una 
    \textbf{congiunzione} ($\land$) di \textbf{clausole}.
\end{definition}

Per capire meglio questi concetti, consideriamo l'esempio seguente:

\begin{exa}
    L'espressione $x_1 \lor \neg x_2 \lor x_3$ è una \textbf{clausola} composta da tre 
    \textit{literals}. 

    Considerando altre due clausole arbitrarie, otteniamo (per esempio) la formula seguente, 
    che è in \textbf{forma CNF}:
    \[
        (x_1 \lor \neg x_2 \lor x_3) \land (\neg x_1 \lor x_4) \land (x_2 \lor x_3)
    \]
\end{exa}

Date queste definizioni, possiamo ora formalizzare il problema classico della soddisfacibilità booleana:

\begin{definition}[PROBLEMA SAT]
    Data una formula logica in \textbf{forma CNF}, il \textbf{problema della 
    soddisfacibilità booleana} (\textbf{SAT}) consiste nel determinare se esiste 
    un \textbf{assegnamento di verità} alle proposizioni atomiche tale che la 
    \textbf{formula risulti vera} (ovvero che \textit{tutte le clausole} siano 
    contemporaneamente soddisfatte).
\end{definition}

Le clausole godono di alcune importanti \textbf{proprietà strutturali} che ci permettono di
giustificare l'uso esclusivo della forma CNF (invece che di altre forme):

\begin{remark}[PROPRIETÀ DELLE CLAUSOLE]
    Notiamo due fatti fondamentali:
    \begin{enumerate}
        \item \textbf{Assorbimento}: una clausola $F$ \textbf{implica} un'altra clausola $G$ 
            se e solo se la \textbf{assorbe}, cioè se \textit{ogni literal} di $F$ è anche 
            un \textit{literal} di $G$. Pertanto, due clausole sono \textbf{equivalenti} 
            se e solo se sono \textbf{uguali}.
        \item \textbf{Generalità della forma CNF}: \textbf{ogni formula logica} può 
            essere trasformata in una formula equivalente in \textbf{forma CNF} e tale 
            trasformazione ha \textbf{complessità polinomiale}. Possiamo quindi considerare 
            formule in CNF \textit{senza perdita di generalità}.
    \end{enumerate}
\end{remark}

\subsection{Clausole di Horn}

Una particolare sottoclasse di clausole riveste un ruolo centrale sia dal punto di vista 
modellistico che computazionale:

\begin{definition}[CLAUSOLA DI HORN]
    Una clausola si dice \textbf{clausola di Horn} se contiene \textbf{al più un literal 
    positivo} (cioè non negato):
    \[
        x_k \lor \left(\bigvee_j \neg x_j\right)
    \]
\end{definition}

Possiamo dare la seguente interpretazione a queste clausole

\begin{remark}[INTERPRETAZIONE E APPLICAZIONI]
    Applicando le equivalenze logiche viste in precedenza, una clausola di Horn può essere 
    riscritta in forma di \textbf{implicazione}:
    \[
        x_k \lor \left(\bigvee_j \neg x_j\right) \equiv \left(\bigwedge_j x_j\right) \to x_k
    \]
    Grazie a questa struttura naturale di regola \textit{"se... allora..."}, le clausole 
    di Horn sono ampiamente utilizzate per \textbf{modellare sistemi esperti}.
\end{remark}

\section{Algoritmi risolutivi per SAT}

Esistono due famiglie principali di algoritmi per risolvere problemi SAT:
\begin{enumerate}
    \item La regola di inferenza di \textbf{risoluzione} (\textit{resolution}).
    \item Un algoritmo \textbf{enumerativo ad albero}, detto \textbf{schema DPLL} 
        (dagli inventori \textit{Davis-Putnam-Logemann-Loveland}).
\end{enumerate}

E' importante sottolinare che, essendo il problema SAT \textbf{NP-completo}, questi 
algoritmi risolutivi presentano una \textbf{complessità esponenziale al caso pessimo}.

Detto ciò, vediamoli più nel dettaglio di seguito.

\subsection{Resolution}

La \textbf{risoluzione} è una regola di inferenza logica che permette di derivare una 
nuova clausola a partire da due clausole esistenti che contengono una coppia di 
\textit{literals} complementari.

\begin{definition}[REGOLA DI RISOLUZIONE]
    Date due \textbf{clausole generatrici} tali che una contenga un literal $x_i$ e l'altra 
    la sua negazione $\neg x_i$:
    \[
        C_1 = (x_i \lor A) \text{ , } \qquad C_2 = (\neg x_i \lor B)
    \]
    La \textbf{regola di risoluzione} permette di inferire una nuova clausola, detta 
    \textbf{risolvente} (\textit{resolvent}), ottenuta unendo i restanti literals:
    \[
        R = (A \lor B)
    \]
\end{definition}

\begin{exa}
    Consideriamo le due clausole generatrici seguenti:
    \[
        \begin{aligned}
            C_1 &: x_1 \lor x_2 \lor x_3 \\
            C_2 &: \neg x_1 \lor x_2 \lor \neg x_4
        \end{aligned}
    \]
    Applicando la regola di risoluzione rispetto alla variabile $x_1$, otteniamo il 
    \textbf{risolvente}:
    \[
        R : x_2 \lor x_3 \lor \neg x_4
    \]
    Come mostrato nella figura~\ref{fig:06_01_RegolaRisoluzione}, i due literals 
    complementari si elidono.
\end{exa}

\begin{figure}[!ht]
    \centering
    \begin{tikzpicture}[>=stealth, scale=1.1]
        % Stile box clausole
        \tikzstyle{clausebox} = [rectangle, draw=nord3, thick, rounded corners, 
        fill=nord6, inner sep=8pt, font=\large]
        \tikzstyle{resbox} = [rectangle, draw=nord14, ultra thick, rounded corners, 
        fill=nord14!15, inner sep=8pt, font=\large]

        % Clausole generatrici
        \node[clausebox] (C1) at (-2.8, 1.8) 
        {$C_1 : \;\; \textcolor{nord10}{\mathbf{x_1}} \lor x_2 \lor x_3$};
        \node[clausebox] (C2) at (2.8, 1.8) 
        {$C_2 : \;\; \textcolor{nord11}{\mathbf{\neg x_1}} \lor x_2 \lor \neg x_4$};

        % Risolvente
        \node[resbox] (R) at (0, -0.3) 
        {$R : \;\; x_2 \lor x_3 \lor \neg x_4$};

        % Frecce di inferenza
        \draw[->, ultra thick, nord3] (C1.south) -- (R.north west);
        \draw[->, ultra thick, nord3] (C2.south) -- (R.north east);
    \end{tikzpicture}
    \caption{Applicazione della regola di inferenza di risoluzione}
    \label{fig:06_01_RegolaRisoluzione}
\end{figure}

L'intuizione alla base della regola è un semplice \textbf{ragionamento per casi}: 
la proposizione atomica $x_1$ può essere solo \textit{vera} oppure \textit{falsa}.
\begin{itemize}
    \item Se $x_1$ è \textbf{vera}, affinché $C_2$ sia soddisfatta deve essere vera 
        $x_2 \lor \neg x_4$.
    \item Se $x_1$ è \textbf{falsa}, affinché $C_1$ sia soddisfatta deve essere vera 
        $x_2 \lor x_3$.
\end{itemize}
In \textbf{entrambi i casi}, il risolvente $R = x_2 \lor x_3 \lor \neg x_4$ risulta 
\textbf{necessariamente vero}.

Detto questo, possiamo descrivere un \textbf{algoritmo risolutivo} per il problema SAT 
basato sulla regola di risoluzione:

\begin{definition}[ALGORITMO DI RISOLUZIONE]
    Dato un insieme $S$ di clausole di partenza, poniamo $S' = S$ e procediamo come segue:
    \begin{enumerate}
        \item \textbf{Finché} $S'$ contiene due clausole con un \textbf{risolvente} $R$ 
            \textit{non implicato} da alcuna clausola già presente in $S'$:
            \begin{enumerate}
                \item \textbf{Rimuoviamo} da $S'$ tutte le clausole \textit{implicate} da $R$ 
                    (per assorbimento).
                \item \textbf{Aggiungiamo} il nuovo risolvente all'insieme: $S' \leftarrow S' \cup \{R\}$.
            \end{enumerate}
        \item Al termine (quando non è più possibile derivare nulla di nuovo), l'insieme 
            $S$ è \textbf{soddisfacibile se e solo se} $S'$ \textbf{non contiene la clausola vuota}.
    \end{enumerate}
\end{definition}

\begin{remark}[ANALOGIA CON CUTTING PLANES]
    L'algoritmo di risoluzione è un \textbf{metodo di inferenza puro} (non effettua alcuna 
    enumerazione) e rappresenta l'\textbf{equivalente in ambito SAT} dell'approccio a 
    \textbf{piani di taglio puri} (\textit{cutting planes}) visto per i problemi MIP.

    Proprio come \textit{cutting planes}, \textbf{non} costituisce lo \textbf{stato dell'arte} per la 
    risoluzione pratica dei problemi SAT.
\end{remark}

\subsubsection{Unit Resolution}

Nella \textit{pratica}, i risolutori SAT implementano una \textbf{"versione ristretta"} ma 
estremamente efficiente della regola di risoluzione:

\begin{definition}[UNIT RESOLUTION]
    Si parla di \textbf{unit resolution} (\textit{risoluzione unitaria}) quando 
    \textbf{almeno una delle due clausole generatrici} è una \textbf{clausola unitaria}, 
    ovvero una clausola costituita da un \textbf{solo literal} (es. $x_1$ oppure $\neg x_1$).
\end{definition}

Questa restrizione comporta due enormi \textbf{vantaggi computazionali}:
\begin{enumerate}
    \item \textbf{Complessità lineare}: può essere eseguita in tempo \textit{lineare} nel 
        numero di clausole (sfruttando strutture dati efficienti come i \textit{watch literals}).
    \item \textbf{Nessuna crescita del problema}: il risolvente $R$ implica sempre la 
        clausola generatrice non unitaria (che viene quindi rimossa per assorbimento), 
        pertanto il \textbf{numero complessivo di clausole non aumenta mai}.
\end{enumerate}

Tuttavia, questa efficienza ha un prezzo dal punto di vista teorico:

\begin{remark}[INCOMPLETEZZA DI UNIT RESOLUTION]
    A differenza della risoluzione generale, \textbf{unit resolution non è un metodo di 
    inferenza completo}: potrebbe infatti \textit{non derivare la clausola vuota} anche 
    quando il sistema di partenza è \textbf{insoddisfacibile}.
\end{remark}

\begin{exa}
    Consideriamo il seguente sistema di quattro clausole sulle variabili $x_1, x_2$:
    \[
        \begin{aligned}
            &x_1 \lor x_2 \\
            &\neg x_1 \lor x_2 \\
            &x_1 \lor \neg x_2 \\
            &\neg x_1 \lor \neg x_2
        \end{aligned}
    \]
    Il sistema è chiaramente \textbf{insoddisfacibile} (nessuna delle quattro combinazioni 
    di verità soddisfa tutte le clausole), ma \textbf{unit resolution non deriva nulla}, 
    in quanto \textit{nessuna clausola è unitaria}.
\end{exa}

Esistono però classi strutturate di problemi per cui \textit{unit resolution} torna ad 
essere \textbf{completa}: 

\begin{prop}[COMPLETEZZA SULLE CLAUSOLE DI HORN]
    Un insieme di \textbf{clausole di Horn} è \textbf{insoddisfacibile se e solo se} 
    \textit{unit resolution} genera la \textbf{clausola vuota}.
\end{prop}

Per capire meglio questo fatto, osserviamo che una \textit{clausola di Horn} può assumere
soltanto \textbf{tre forme}:
\begin{enumerate}
    \item \textbf{Clausola unitaria}: $x_k$ oppure $\neg x_k$.
    \item \textbf{Implicazione non unitaria}: $\neg x_{1} \lor \ldots \lor \neg x_m \lor x_k$,
        con $m\ge 1$ e quindi lunghezza $\ge 2$.
    \item \textbf{Puramente negativa non unitaria}: $\neg x_{1} \lor \ldots \lor \neg x_m$
        con $m\ge 2$.
\end{enumerate}

Da questa osservazione derivano due \textbf{situazioni possibili}:
\begin{enumerate}
    \item Se nell'insieme \textbf{non} ci sono \textbf{clausole unitarie}, allora ogni
        clausola rientra necessariamente nei \textbf{casi 2 o 3} e quindi
        \textbf{contiene necessariamanete almeno un literal negato}.
        Ma allora l'assegnamento banale in cui \textbf{poniamo tutte le variabili a falso}
        rende \textbf{veri tutti i literals negati} e \textbf{soddisfa l'intera formula}.
    \item Se con l'\textbf{assegnamento banale} sopra l'insieme è \textbf{insoddisfacibile},
        allora doveva esserci \textit{almeno una clausola contentente solo literal non negati},
        che devi quindi essere del \textbf{tipo 1}.
\end{enumerate}

\begin{remark}
    Fatte queste osservazioni, è immediato capire che:
    \begin{enumerate}
        \item Nel primo caso, \textbf{unit resolution \textit{non fa niente}}, per cui
            \textbf{non genera la clausola vuota}, ed il problema è \textbf{soddisfacibile}.
        \item Nel secondo caso, \textbf{esiste una clausola unitaria} e quindi
            \textbf{unit resolution può procedere}.
    \end{enumerate}
\end{remark}

\subsection{Algoritmo DPLL}

Così come per i problemi MIP abbiamo combinato piani di taglio e ricerca ad albero nel 
\textit{Branch and Cut}, lo stato dell'arte per la risoluzione di problemi SAT si basa 
sullo \textbf{schema DPLL} (\textit{Davis-Putnam-Logemann-Loveland}), un 
\textbf{algoritmo enumerativo di ricerca ad albero}.

Vediamolo descritto di seguito:

\begin{definition}[SCHEMA DPLL]
    L'algoritmo esplora uno spazio di ricerca ad albero operando come segue:
    \begin{enumerate}
        \item \textbf{Inferenza (semplificazione del nodo)}: ad ogni nodo dell'albero si 
            applicano fino a convergenza due sole tecniche di inferenza:
            \begin{enumerate}
                \item \textbf{Unit resolution}: propaga i vincoli forzati dalle clausole unitarie.
                \item \textbf{Pure literal elimination}: fissa direttamente il valore di verità 
                    dei \textit{literals} che compaiono in tutte le clausole attive con lo 
                    \textit{stesso segno} (sempre positivi o sempre negati).
            \end{enumerate}
        \item \textbf{Verifica di terminazione}:
            \begin{enumerate}
                \item Se \textbf{tutte le clausole sono soddisfatte}, abbiamo trovato un 
                    assegnamento valido e \textbf{terminiamo}.
                \item Se viene derivata la \textbf{clausola vuota}, il ramo 
                    corrente è \textbf{inconsistente} e si effettua \textbf{backtracking}.
            \end{enumerate}
        \item \textbf{Branching}: se il problema non è ancora deciso, si sceglie una 
            variabile $x_i$ non ancora assegnata e si generano \textbf{due sottoproblemi} 
            fissandone il valore di verità (ovvero aggiungendo la clausola unitaria $x_i$ 
            nel ramo sinistro e $\neg x_i$ nel ramo destro).
    \end{enumerate}
\end{definition}

L'operazione di \textit{branching} dello schema DPLL (mostrata nella 
figura~\ref{fig:06_02_DPLL_BranchTree}) è del tutto analoga al \textit{branching} 
sulle variabili binarie visto nel Capitolo 4: aggiungere una clausola unitaria rende 
inoltre \textit{unit resolution} immediatamente applicabile nel sottoproblema figlio.

\begin{figure}[!ht]
    \centering
    \begin{tikzpicture}[scale = 0.8]
        % Stile dei nodi dell'albero (coerente con Cap. 4)
        \tikzstyle{treenode} = [circle, draw=nord3, ultra thick, fill=nord6, minimum size=1.2cm]

        % Nodo Radice (Formula corrente)
        \node[treenode] (root) at (0, 0) {};

        % Nodi Figli (Sottoproblemi con clausola unitaria aggiunta)
        \node[treenode] (left) at (-1.8, -1.5) {};
        \node[treenode] (right) at (1.8, -1.5) {};

        % Rami (Operazione di Branching)
        \draw[ultra thick, nord3] (root) -- (left);
        \draw[ultra thick, nord3] (root) -- (right);

        % Etichette dei vincoli aggiunti
        \node[below=0.2cm, font=\large\bfseries, nord10, align=center] at (left.south) 
        {$x_i = T$ \\ \footnotesize\textcolor{nord3}{(aggiungo $x_i$)}};
        \node[below=0.2cm, font=\large\bfseries, nord11, align=center] at (right.south) 
        {$x_i = F$ \\ \footnotesize\textcolor{nord3}{(aggiungo $\neg x_i$)}};
    \end{tikzpicture}
    \caption{Operazione di branching nello schema DPLL}
    \label{fig:06_02_DPLL_BranchTree}
\end{figure}

\begin{remark}[GESTIONE DELLA MEMORIA E SOLVER MODERNI]
    Concludiamo con due osservazioni fondamentali sulle implementazioni pratiche:
    \begin{itemize}
        \item \textbf{Occupazione di memoria lineare}: l'albero di ricerca DPLL viene 
            visitato in \textbf{Depth-First Search (DFS)}, il che garantisce un'occupazione 
            di memoria \textit{lineare} nella dimensione del problema.
        \item \textbf{Estensioni stato dell'arte}: i moderni risolutori SAT (\textit{CDCL}, 
            Conflict-Driven Clause Learning) potenziano lo schema DPLL con tecniche 
            avanzate quali il \textbf{conflict learning} (apprendimento di nuove clausole 
            dai conflitti incontrati per potare rami futuri) e i \textbf{restart} periodici.
    \end{itemize}
\end{remark}
